% ============================================================================
% DM-2026-011: Claude Managed Agents Adoption for the Security Toolkit
%   Decision 1 — Agent runtime surface for the scan-and-attest agent
%   Decision 2 — Trust boundary: sandbox type by target sensitivity
%   Decision 3 — Scope across the other repositories
% ============================================================================
\newcommand{\UniqueID}{DM-2026-011}
\newcommand{\DocumentDate}{September 13, 2026}
\newcommand{\AuthorName}{Bruce Dombrowski}
\newcommand{\AuthorTitle}{Systems Engineer, Payload Software Integration}
\newcommand{\ToField}{Security Toolkit Stakeholders; AI-Agents Repository Maintainers}
\newcommand{\SubjectField}{Adoption of Claude Managed Agents for the Security Toolkit Scan-and-Attest Agent}
\newcommand{\OPRField}{Payload Software Integration (PSI)}

\newcommand{\dmContent}{%
\begin{enumerate}

\dmsection{Purpose}

This memorandum records the decision to adopt Anthropic's Claude Managed Agents
platform (public beta, April 2026) as the runtime for a scan-and-attest agent
built on the Security Toolkit, the trust-boundary rules that govern where the
agent's tools may execute, and the scope of adoption across the maintainer's
other repositories. It follows the pattern presented in the Code with Claude
2026 session \textit{Ship your first Managed Agent} (I. He, Anthropic Applied
AI), whose reference implementation is published in
\texttt{anthropics/cwc-workshops}.

\dmsection{Background}

Managed Agents separates the agent loop (``brain'', run on Anthropic's
orchestration layer) from tool execution (``hands'', run in a per-session
sandbox). Three persisted resources define a run: an \textbf{Agent} (model,
system prompt, tools, skills; versioned, created once), an
\textbf{Environment} (sandbox template: cloud or self-hosted), and a
\textbf{Session} (one run, with an event log). Custom tools execute on the
caller's host, so secrets and privileged actions never enter the sandbox.
Skills in a mounted repository's \texttt{.claude/skills} directory are loaded
at session start. Sessions accept a hard dollar budget and persist a complete
event history.

The Security Toolkit already carries the pieces such an agent needs: a
non-interactive scan entry point (\texttt{run-all-scans.sh -n}), attestation
generation with toolkit version and commit hashes, an allowlist model that
requires a SHA256 integrity hash plus a justification, and TIA principles
(Transparency, Inspectability, Accountability, Traceability) that map directly
onto a persisted session event log.

\dmsection{Decision 1 --- Agent Runtime Surface}

\textbf{Options considered:}
\begin{itemize}
    \item \textbf{1A --- Claude API with tool runner, self-hosted loop.} Maximum
    control; the maintainer hosts the loop, sandbox, scheduler, and state.
    \item \textbf{1B --- Claude Agent SDK.} Claude Code packaged as a library;
    batteries included, still self-hosted.
    \item \textbf{1C --- Claude Managed Agents (SELECTED).} Anthropic runs the loop
    and (optionally) the sandbox; agent configuration is a version-controlled
    YAML manifest applied once with the \texttt{ant} CLI; sessions are created by
    a thin runner script.
    \item \textbf{1D --- No agent.} Continue running the toolkit by hand and in CI.
\end{itemize}

\textbf{Decision:} Option 1C. The agent definition lives in the toolkit
repository as \texttt{agents/security-attest.agent.yaml} and
\texttt{agents/security-scan.env.yaml}; the runbook lives in
\texttt{.claude/skills/security-scan-runbook/SKILL.md}; the runner is
\texttt{agents/run\_scan\_session.py}. Two host-side custom tools are defined:
\texttt{get\_kev\_status} (CISA KEV lookup from the host cache) and
\texttt{request\_allowlist\_entry} (a request only; the human decides and the
host computes the integrity hash).

\textbf{Rationale:} The toolkit's value is auditability, and Managed Agents
gives a persisted, inspectable event log per session at no implementation
cost. Versioned agent objects match the toolkit's own attestation discipline
(every attestation names a toolkit version and commit; every session pins an
agent version). Options 1A and 1B require the maintainer to build and host the
loop, sandbox, and state store for a solo-maintained project. Option 1D
forgoes the judgment step (classifying findings as true or false positives)
that consumes most of the human time today.

\dmsection{Decision 2 --- Trust Boundary by Target Sensitivity}

\textbf{Context:} A cloud environment clones the target repository into an
Anthropic-hosted container. The toolkit is public and may be mounted anywhere.
Targets are not.

\textbf{Decision:}
\begin{itemize}
    \item Targets that are public or contain no CUI and no PII \textbf{may} use the
    cloud environment.
    \item Targets that contain or may contain CUI, PII, or export-controlled data
    \textbf{shall} use a \texttt{self\_hosted} environment, where tool execution
    and file contents remain on maintainer-controlled infrastructure. Web search
    and web fetch are disabled on the agent regardless of environment.
    \item The agent \textbf{shall not} modify allowlists, delete files, or run
    scripts from the target repository. Target contents are treated as untrusted
    data; instruction-like text found there is a finding, not a command.
    \item Only trusted repositories are mounted, because a mounted repository's
    \texttt{.claude/skills} directory is loaded as agent instructions without
    review.
    \item Every session is created with a dollar budget cap.
\end{itemize}

\textbf{Rationale:} These rules keep the CUI determination (DM-2026-007,
DM-2026-008) intact: nothing the desk instruction classifies as CUI leaves
maintainer-controlled systems. The allowlist rule preserves Accountability;
the untrusted-data rule addresses prompt injection through scanned content,
which is the primary new risk an agent introduces to a scanner.

\dmsection{Decision 3 --- Scope Across Other Repositories}

\begin{itemize}
    \item \textbf{Security Toolkit:} adopt now (this memorandum).
    \item \textbf{Foreclosures:} candidate for a monthly scheduled deployment that
    runs the deterministic pipeline and reviews the ranked shortlist. Deferred
    until the toolkit agent has run cleanly for one release cycle.
    \item \textbf{Investing:} a read-only daily review agent over the run report and
    decision log is acceptable; brokerage credentials go in a vault with
    read-only scope. The agent \textbf{shall never} be given order submission.
    Deferred.
    \item \textbf{Email Assistant:} excluded. Its stated requirement is that no
    content leaves the laptop and no cloud model is used; Managed Agents
    contradicts that by design.
    \item \textbf{Philosophy:} excluded; no identified benefit.
\end{itemize}

\dmsection{Implementation and Constraints}

\begin{itemize}
    \item \texttt{.gitignore} changed from \texttt{.claude/} to
    \texttt{.claude/*} with \texttt{!.claude/skills/} so the runbook is committed
    and discoverable when the toolkit is mounted; local Claude Code settings
    remain ignored.
    \item Agent and environment are created once with \texttt{agents/setup.sh};
    the returned IDs are stored in \texttt{agents/ids.env} (git-ignored). Sessions
    reference IDs only; no code path calls \texttt{agents.create()} per run.
    \item Model: \texttt{claude-opus-5} at high effort. Tools: bash, read, grep,
    glob, write (for the report), web tools disabled.
    \item Dependencies: \texttt{ant} CLI and the \texttt{anthropic} Python package
    on the host; both are optional for the toolkit's existing Bash workflow.
    \item CHANGELOG entry under Unreleased; regression test
    \texttt{tests/test-managed-agent-manifests.sh} validates manifests, runner
    syntax, and that the skill directory is not git-ignored.
\end{itemize}

\dmsection{Verification}

\begin{itemize}
    \item \textbf{Manifest test:} both YAML manifests parse; the runner compiles;
    \texttt{git check-ignore} reports the skill file as not ignored.
    \item \textbf{Dry session:} against the toolkit's own \texttt{demo/} planted
    findings (\texttt{docs/DEMO-PLANTED-FINDINGS.md}), the agent reports every
    planted finding with a correct \texttt{path:line} and requests, but does not
    perform, allowlist changes.
    \item \textbf{Boundary test:} a target containing instruction-like text is
    reported as a finding and the instruction is not followed.
    \item \textbf{Budget test:} a session created with a budget pauses with
    \texttt{budget\_reached} rather than running past the cap.
\end{itemize}

\dmsection{References}

\begin{itemize}
    \item Code with Claude 2026, \textit{Ship your first Managed Agent}, I. He;
    reference code \texttt{anthropics/cwc-workshops/ship-your-first-managed-agent}.
    \item Claude Managed Agents documentation (platform.claude.com), beta
    \texttt{managed-agents-2026-04-01}.
    \item Security Toolkit \texttt{CLAUDE.md} (TIA principles), \texttt{AGENTS.md},
    \texttt{docs/COMPLIANCE.md}.
    \item DM-2026-007, DM-2026-008 (CUI determination and desk instruction).
    \item NIST SP 800-53 Rev.\ 5, NIST SP 800-171, FIPS 199/200; RFC 2119.
\end{itemize}

\end{enumerate}
}

% ============================================================================
% Decision Memorandum Base Template
% ============================================================================
% This is the shared base template for all Decision Memoranda.
% DO NOT edit this file for individual decisions.
%
% Usage: Create a thin wrapper that defines variables and content, then:
%   % ============================================================================
% Decision Memorandum Base Template
% ============================================================================
% This is the shared base template for all Decision Memoranda.
% DO NOT edit this file for individual decisions.
%
% Usage: Create a thin wrapper that defines variables and content, then:
%   % ============================================================================
% Decision Memorandum Base Template
% ============================================================================
% This is the shared base template for all Decision Memoranda.
% DO NOT edit this file for individual decisions.
%
% Usage: Create a thin wrapper that defines variables and content, then:
%   \input{_template.tex}
%
% Required variables (define before \input):
%   \UniqueID       - e.g., DM-2026-001
%   \DocumentDate   - e.g., January 19, 2026
%   \AuthorName     - Signer name
%   \AuthorTitle    - Signer title
%   \ToField        - Distribution
%   \SubjectField   - Subject line
%   \OPRField       - Office of Primary Responsibility
%   \dmContent      - The full content (enumerate environment with \dmsection items)
%
% ============================================================================

\documentclass[11pt,letterpaper]{article}

% ============================================================================
% PACKAGES
% ============================================================================
\usepackage[margin=1in, top=1.15in, headheight=30pt]{geometry}
\usepackage{graphicx}
\usepackage{fancyhdr}
\usepackage{lastpage}
\usepackage{enumitem}
\usepackage{xcolor}
\usepackage{hyperref}
\usepackage{tabularx}

% ============================================================================
% DOCUMENT CONFIGURATION
% ============================================================================
\hypersetup{
    colorlinks=true,
    linkcolor=blue,
    urlcolor=blue,
    pdfborder={0 0 0}
}

% Remove paragraph indentation
\setlength{\parindent}{0pt}
\setlength{\parskip}{6pt}

% List formatting - match Word style
\setlist[enumerate]{leftmargin=0.5in, labelwidth=0.25in, labelsep=0.1in, align=left, itemsep=12pt, topsep=12pt}
\setlist[enumerate,1]{label=\arabic*.}
\setlist[itemize]{leftmargin=0.5in, itemsep=3pt, topsep=6pt}

% ============================================================================
% HEADER AND FOOTER CONFIGURATION
% ============================================================================
\pagestyle{fancy}
\fancyhf{}

% Header: text-only, Boeing / ISS Program identification
\fancyhead[L]{\large\textbf{Decision Memorandum}\\[-2pt]\scriptsize Boeing \textbar{} International Space Station Program}
\fancyhead[R]{\small \UniqueID\\[-2pt]\scriptsize \DocumentDate}
\renewcommand{\headrulewidth}{0.4pt}

% Footer: Unique ID (left), Page X of Y (center), Date (right)
\fancyfoot[L]{\small \UniqueID}
\fancyfoot[C]{\small Page \thepage\ of \pageref{LastPage}}
\fancyfoot[R]{\small \DocumentDate}
\renewcommand{\footrulewidth}{0.4pt}

% ============================================================================
% CUSTOM COMMANDS
% ============================================================================
% Bold section headers for numbered sections
\newcommand{\dmsection}[1]{\item \textbf{#1}\par}

% ============================================================================
% BEGIN DOCUMENT
% ============================================================================
\begin{document}

% ----------------------------------------------------------------------------
% UNIQUE ID (top right)
% ----------------------------------------------------------------------------
\hfill \UniqueID

\vspace{0.25in}

% ----------------------------------------------------------------------------
% SIGNATURE BLOCK
% ----------------------------------------------------------------------------
\underline{\hspace{2.5in}}\\[2pt]
\AuthorName\\
\AuthorTitle

\vspace{0.25in}

% ----------------------------------------------------------------------------
% MEMO HEADER FIELDS
% ----------------------------------------------------------------------------
\textbf{DATE:} \DocumentDate

\textbf{TO:} \ToField

\textbf{SUBJECT:} \SubjectField

\textbf{OFFICE OF PRIMARY RESPONSIBILITY:} \OPRField

\vspace{0.25in}

% ============================================================================
% MAIN CONTENT - NUMBERED SECTIONS
% ============================================================================
\dmContent

\end{document}

%
% Required variables (define before \input):
%   \UniqueID       - e.g., DM-2026-001
%   \DocumentDate   - e.g., January 19, 2026
%   \AuthorName     - Signer name
%   \AuthorTitle    - Signer title
%   \ToField        - Distribution
%   \SubjectField   - Subject line
%   \OPRField       - Office of Primary Responsibility
%   \dmContent      - The full content (enumerate environment with \dmsection items)
%
% ============================================================================

\documentclass[11pt,letterpaper]{article}

% ============================================================================
% PACKAGES
% ============================================================================
\usepackage[margin=1in, top=1.15in, headheight=30pt]{geometry}
\usepackage{graphicx}
\usepackage{fancyhdr}
\usepackage{lastpage}
\usepackage{enumitem}
\usepackage{xcolor}
\usepackage{hyperref}
\usepackage{tabularx}

% ============================================================================
% DOCUMENT CONFIGURATION
% ============================================================================
\hypersetup{
    colorlinks=true,
    linkcolor=blue,
    urlcolor=blue,
    pdfborder={0 0 0}
}

% Remove paragraph indentation
\setlength{\parindent}{0pt}
\setlength{\parskip}{6pt}

% List formatting - match Word style
\setlist[enumerate]{leftmargin=0.5in, labelwidth=0.25in, labelsep=0.1in, align=left, itemsep=12pt, topsep=12pt}
\setlist[enumerate,1]{label=\arabic*.}
\setlist[itemize]{leftmargin=0.5in, itemsep=3pt, topsep=6pt}

% ============================================================================
% HEADER AND FOOTER CONFIGURATION
% ============================================================================
\pagestyle{fancy}
\fancyhf{}

% Header: text-only, Boeing / ISS Program identification
\fancyhead[L]{\large\textbf{Decision Memorandum}\\[-2pt]\scriptsize Boeing \textbar{} International Space Station Program}
\fancyhead[R]{\small \UniqueID\\[-2pt]\scriptsize \DocumentDate}
\renewcommand{\headrulewidth}{0.4pt}

% Footer: Unique ID (left), Page X of Y (center), Date (right)
\fancyfoot[L]{\small \UniqueID}
\fancyfoot[C]{\small Page \thepage\ of \pageref{LastPage}}
\fancyfoot[R]{\small \DocumentDate}
\renewcommand{\footrulewidth}{0.4pt}

% ============================================================================
% CUSTOM COMMANDS
% ============================================================================
% Bold section headers for numbered sections
\newcommand{\dmsection}[1]{\item \textbf{#1}\par}

% ============================================================================
% BEGIN DOCUMENT
% ============================================================================
\begin{document}

% ----------------------------------------------------------------------------
% UNIQUE ID (top right)
% ----------------------------------------------------------------------------
\hfill \UniqueID

\vspace{0.25in}

% ----------------------------------------------------------------------------
% SIGNATURE BLOCK
% ----------------------------------------------------------------------------
\underline{\hspace{2.5in}}\\[2pt]
\AuthorName\\
\AuthorTitle

\vspace{0.25in}

% ----------------------------------------------------------------------------
% MEMO HEADER FIELDS
% ----------------------------------------------------------------------------
\textbf{DATE:} \DocumentDate

\textbf{TO:} \ToField

\textbf{SUBJECT:} \SubjectField

\textbf{OFFICE OF PRIMARY RESPONSIBILITY:} \OPRField

\vspace{0.25in}

% ============================================================================
% MAIN CONTENT - NUMBERED SECTIONS
% ============================================================================
\dmContent

\end{document}

%
% Required variables (define before \input):
%   \UniqueID       - e.g., DM-2026-001
%   \DocumentDate   - e.g., January 19, 2026
%   \AuthorName     - Signer name
%   \AuthorTitle    - Signer title
%   \ToField        - Distribution
%   \SubjectField   - Subject line
%   \OPRField       - Office of Primary Responsibility
%   \dmContent      - The full content (enumerate environment with \dmsection items)
%
% ============================================================================

\documentclass[11pt,letterpaper]{article}

% ============================================================================
% PACKAGES
% ============================================================================
\usepackage[margin=1in, top=1.15in, headheight=30pt]{geometry}
\usepackage{graphicx}
\usepackage{fancyhdr}
\usepackage{lastpage}
\usepackage{enumitem}
\usepackage{xcolor}
\usepackage{hyperref}
\usepackage{tabularx}

% ============================================================================
% DOCUMENT CONFIGURATION
% ============================================================================
\hypersetup{
    colorlinks=true,
    linkcolor=blue,
    urlcolor=blue,
    pdfborder={0 0 0}
}

% Remove paragraph indentation
\setlength{\parindent}{0pt}
\setlength{\parskip}{6pt}

% List formatting - match Word style
\setlist[enumerate]{leftmargin=0.5in, labelwidth=0.25in, labelsep=0.1in, align=left, itemsep=12pt, topsep=12pt}
\setlist[enumerate,1]{label=\arabic*.}
\setlist[itemize]{leftmargin=0.5in, itemsep=3pt, topsep=6pt}

% ============================================================================
% HEADER AND FOOTER CONFIGURATION
% ============================================================================
\pagestyle{fancy}
\fancyhf{}

% Header: text-only, Boeing / ISS Program identification
\fancyhead[L]{\large\textbf{Decision Memorandum}\\[-2pt]\scriptsize Boeing \textbar{} International Space Station Program}
\fancyhead[R]{\small \UniqueID\\[-2pt]\scriptsize \DocumentDate}
\renewcommand{\headrulewidth}{0.4pt}

% Footer: Unique ID (left), Page X of Y (center), Date (right)
\fancyfoot[L]{\small \UniqueID}
\fancyfoot[C]{\small Page \thepage\ of \pageref{LastPage}}
\fancyfoot[R]{\small \DocumentDate}
\renewcommand{\footrulewidth}{0.4pt}

% ============================================================================
% CUSTOM COMMANDS
% ============================================================================
% Bold section headers for numbered sections
\newcommand{\dmsection}[1]{\item \textbf{#1}\par}

% ============================================================================
% BEGIN DOCUMENT
% ============================================================================
\begin{document}

% ----------------------------------------------------------------------------
% UNIQUE ID (top right)
% ----------------------------------------------------------------------------
\hfill \UniqueID

\vspace{0.25in}

% ----------------------------------------------------------------------------
% SIGNATURE BLOCK
% ----------------------------------------------------------------------------
\underline{\hspace{2.5in}}\\[2pt]
\AuthorName\\
\AuthorTitle

\vspace{0.25in}

% ----------------------------------------------------------------------------
% MEMO HEADER FIELDS
% ----------------------------------------------------------------------------
\textbf{DATE:} \DocumentDate

\textbf{TO:} \ToField

\textbf{SUBJECT:} \SubjectField

\textbf{OFFICE OF PRIMARY RESPONSIBILITY:} \OPRField

\vspace{0.25in}

% ============================================================================
% MAIN CONTENT - NUMBERED SECTIONS
% ============================================================================
\dmContent

\end{document}

